\documentclass[10pt,a4paper]{amsart}
\usepackage[margin=19mm]{geometry}
\usepackage{amsmath,amssymb,mathtools}
\usepackage[scr=rsfs,cal=euler]{mathalfa}
\usepackage{xfrac}
\usepackage[unicode,hidelinks,bookmarks=false]{hyperref}
\newcommand{\ie}{i.e.\ }
\newcommand{\cf}{cf.\ }
\newcommand{\resp}{respectively\ }
\newcommand{\Subsection}{\subsection}
\providecommand{\nobreakdash}{\nobreak\hbox{-}}
\setlength{\emergencystretch}{2em}
\multlinegap0pt
\usepackage{bm}
\usepackage{bbm}
\usepackage{dsfont}
\usepackage{xspace}
\usepackage{cancel}
\usepackage{mathtools}
\usepackage{amsthm}
\makeatletter
\@ifundefined{theorem}{\newtheorem{theorem}{Theorem}}{}
\@ifundefined{corollary}{\newtheorem{corollary}[theorem]{Corollary}}{}
\@ifundefined{lemma}{\newtheorem{lemma}[theorem]{Lemma}}{}
\@ifundefined{proposition}{\newtheorem{proposition}[theorem]{Proposition}}{}
\makeatother
\DeclareMathOperator{\Hom}{Hom}
\DeclareMathOperator{\Sing}{Sing}
\newcommand{\CAlg}{\mathrm{CAlg}}
\newcommand{\C}{\mathbb{C}}
\newcommand{\Fin}{\mathrm{Fin}}
\newcommand{\Kdef}{\mathcal{K}^{\mathrm{def}}}
\newcommand{\N}{\mathbb{N}}
\newcommand{\Rdef}{\mathcal{R}^{\mathrm{def}}}
\newcommand{\Spc}{\mathrm{Spc}}
\newcommand{\co}{\colon\thinspace}
\title[Spaces of homomorphisms, formality and Hochschild homology]{Spaces of homomorphisms, formality and Hochschild homology}
\author{Simon Gritschacher}
\date{Source: arXiv:2507.17683v1 (CC BY 4.0)}
\begin{document}
\maketitle

\noindent\emph{Source and licence.} Spaces of homomorphisms, formality and Hochschild homology. Version arXiv:2507.17683v1 (CC BY 4.0), distributed under the Creative Commons Attribution 4.0 International licence, as recorded in the arXiv licence metadata for this exact version and shown on the abstract page https://arxiv.org/abs/2507.17683v1. Full rights notice: licenses/arxiv-2507.17683-CC-BY-4.0.txt.\par\smallskip

\noindent\textit{Editorial scope.} Counted unit: the Proposition whose source label is \texttt{prop:key} in the article (source0.tex lines 691--698, source label \texttt{prop:key}), delivered together with the article's own complete original proof of it (source0.tex lines 699--748, one \texttt{proof} environment of 50 lines). No mathematics is written, repaired or re-derived: statement and proof are the article's own text line for line. Required context retained uncounted: eq:equivalencerelation (lines 275--277); lem:singular (lines 465--471); cor:tensor2 (lines 570--576). Every definition, display and label of those lines is retained unchanged. Omitted from this package and not redistributed: 1--274, 278--464, 472--569, 577--690, 749--2048. Nothing inside a retained excerpt is omitted or altered: every retained body is delivered exactly as the article's lines are written, so no omission receipt of any kind accompanies this package. Citations: the retained excerpts carry no citation command at all, so no bibliography entry of the article is transcribed and no reference is redistributed. Reference adaptation: none was needed and none was made; every reference inside the delivered unit resolves inside it, so no unresolved reference or citation is produced. Printed slips kept literal: no printed word, symbol, hypothesis or number of the counted statement or of its proof was altered. Layout and class: the article's own class is \texttt{article} with packages including fullpage, hyperref, amsmath, amssymb, amsthm, graphicx, caption, tikz, amsfonts, dsfont, mathtools, enumerate, eucal, mathrsfs; the wrapper is the lane's pinned \texttt{amsart} editorial wrapper at 10pt on A4, which restates the theorem-like environments the retained text needs and adds the article's own macro definitions that the retained text needs (\texttt{\textbackslash C}, \texttt{\textbackslash CAlg}, \texttt{\textbackslash Fin}, \texttt{\textbackslash Hom}, \texttt{\textbackslash Kdef}, \texttt{\textbackslash N}, \texttt{\textbackslash Rdef}, \texttt{\textbackslash Sing}, \texttt{\textbackslash Spc}, \texttt{\textbackslash co}), with extra wrapper packages (bm, bbm, dsfont, xspace, cancel, mathtools). The delivered numbering is the unit's own. Assets: no retained span contains a figure environment, a graphics-inclusion command, a PDF-page inclusion, a table or a TikZ picture; no image file is copied into this package, the package declares no asset dependency and no image file is redistributed with it. Licence: version arXiv:2507.17683v1 of this article is distributed under the Creative Commons Attribution 4.0 International licence, as recorded in the arXiv licence metadata for this exact version and shown on the abstract page https://arxiv.org/abs/2507.17683v1. A full rights notice is delivered at licenses/arxiv-2507.17683-CC-BY-4.0.txt. Characters: every retained excerpt is pure ASCII, so the delivered engine, XeLaTeX with its default Latin Modern text font, reports no missing character.
\begin{equation} \label{eq:equivalencerelation}
(A(\alpha)(x),f)\sim (x,f \alpha)
\end{equation}

\begin{lemma} \label{lem:singular}
If $A$ is a cofibrant $\Gamma$-space and $X$ is a finite based simplicial set, then there is a natural weak equivalence
\[
(\Sing A)_X \simeq \Sing A_{|X|}\, ,
\]
which is also natural in $A$ and $X$.
\end{lemma}

\begin{corollary} \label{cor:tensor2}
For every space $X\in \Spc$, there is an equivalence in $\CAlg(\Spc^\N)$
\[
\mathrm{N}(A_{X_+})\simeq \mathrm{N}(A)\otimes X
\]
which is natural in $A$ and $X$.
\end{corollary}

\begin{proposition} \label{prop:key}
Let $G,A$ be finitely generated discrete groups and $A$ abelian. Then twisting $G$-representations by characters of $A$ induces equivalences of $\mathbb{E}_\infty$-algebras in $\Spc^\N$
\begin{enumerate}
\item[(1)] $\Kdef(G\times A)\simeq \Kdef(G) \otimes \widehat{A}$
\item[(2)] $\Rdef(G\times A)\simeq \Rdef(G) \otimes \widehat{A}$
\end{enumerate}
and these equivalences are natural in $G$ and $A$.
\end{proposition}

\begin{proof}
We shall only prove (1) as the proof of (2) is analogous.

We will prove an isomorphism of $\Gamma$-spaces
\[
\Kdef(G)_{\widehat{A}_+} \cong \Kdef(G\times A) 
\]
and subsequently observe that this isomorphism respects the graded structure. Thus, it induces an isomorphism of $\N^\otimes$-spaces, and this implies the claim by Lemma \ref{lem:singular} and Corollary \ref{cor:tensor2}.



Let $V\subseteq\C^\infty$ be a finite dimensional linear subspace. A homomorphism $f\co G\times A\to U(V)$ corresponds to a pair $(f',\phi)$ where $f'\co G\to U(V)$ and $\phi \co A\to U(V)$ are homomorphisms such that $f'(g)\phi(a)=\phi(a)f'(g)$ for all $(g,a)\in G\times A$. Since $A$ is abelian, $\phi$ splits $V$ into an orthogonal sum $\bigoplus_{i=1}^k V_i$ such that the action of $\phi$ on each summand is determined by a character $\phi_i\in \Hom(A,U(1))=\widehat{A}$, all $V_i$ are non-zero and the $\phi_i$ are mutually distinct. By commutativity of $f'$ and $\phi$, $f'$ restricts to representations $f_i'\co G\to U(V_i)$. This gives a bijective correspondence
\begin{equation} \label{eq:correspondence}
(V,f)\mapsto \{(V_i,f_i',\phi_i)\}_{i=1}^k\, .
\end{equation}





Fix $\langle l\rangle \in \Fin_\ast$ and $m\geq 0$ and define a homeomorphism
\[
g_{m,\langle l\rangle}\co \Kdef(G)_{\leq m}(\langle l\rangle\wedge \widehat{A}_+) \to \Kdef(G\times A)_{\leq m}(\langle l\rangle)\,.
\]
as follows. Recall that the left hand side is a quotient space of the coproduct
\[
\bigsqcup_{k\geq 0} \Kdef(G)_{\leq m}(\langle k\rangle) \times [(\langle l \rangle^\circ \times \widehat{A})_{+}]^k\,.
\]
Under the equivalence relation (\ref{eq:equivalencerelation}) every element in the coproduct is equivalent to one of the form
\[
((V_i,f'_i)_{i=1}^k,(n_i,\phi_i)_{i=1}^k),
\]
that is, to one where the extra basepoint of $(\langle l\rangle^\circ \times \widehat{A})_+$ does not appear. To such an element we associate
\[
(W_j,f_j)_{j=1}^l \in \Kdef(G\times A)_{\leq m}(\langle l\rangle)
\]
with $(W_j,f_j)$ defined by
\[
(W_j,f_j) \mapsto \{(V_i,f_i',\phi_i)\mid i: n_i=j\}
\]
under the correspondence (\ref{eq:correspondence}). This defines the continuous bijection $g_{m,\langle l\rangle}$; as the domain of $g_{m,\langle l\rangle}$ is compact and the codomain is Hausdorff, $g_{m,\langle l\rangle}$ is in fact a homeomorphism.

For $(x_1,\dots,x_l)\in \N^l$, the element $((V_i,f'_i)_{i=1}^k,(n_i,\phi_i)_{i=1}^k)$ defines an element of the subspace $\Kdef(G)_{\widehat{A}_+}(x_1,\dots,x_l)$ precisely if $\sum_{i: n_i=j} \dim_{\C}(V_i)=x_j$ for all $j=1,\dots,l$. It is therefore mapped to a tuple $(W_j,f_j)_{j=1}^l$ with $\dim_{\C}(W_j)=x_j$, which shows that the isomorphism is in fact one of $\N^\otimes$-spaces.

Taking the colimit over $m$, and noting that $ \Kdef(G)(\langle l\rangle \wedge \widehat{A}_+)$ carries the colimit topology as $\langle l \rangle \wedge \widehat{A}_+$ is compact Hausdorff, we hence obtain a homeomorphism
\[
g_{\langle l\rangle}\co \Kdef(G)(\langle l \rangle\wedge \widehat{A}_+) \to \Kdef(G\times A)(\langle l \rangle)\, .
\]
By inspection, this map is natural in $\langle l \rangle\in \Fin_\ast$ and thus defines an isomorphism of $\N^\otimes$-spaces. Evidently, this isomorphism is natural in $G$ and $A$.
\end{proof}

\end{document}
